\documentclass[11pt, a4paper]{article}

\usepackage[T1]{fontenc}
\usepackage[utf8]{inputenc}
\usepackage{tgpagella}
\usepackage[spanish, es-noshorthands]{babel}
\usepackage{amsmath, amssymb, bm}
\usepackage{booktabs}
\usepackage{tabularx}
\usepackage[table, xcdraw]{xcolor}
\usepackage[most]{tcolorbox}
\usepackage[margin=2.5cm]{geometry}
\usepackage{ragged2e}
\usepackage{fancyhdr}

\pagestyle{fancy}
\fancyhf{}
\setlength{\headheight}{16pt}
\lhead{\textbf{UCB Sede Tarija} | Ing. Mecatrónica}
\chead{}
\rhead{IMT-121 Circuitos Electrónicos I | Tabla de Registro A}
\lfoot{Prof: M.Sc. Ing. Bernardo Quiroga Turdera}
\rfoot{Página \thepage}
\renewcommand{\headrulewidth}{0.4pt}

\definecolor{ucbblue}{RGB}{0, 51, 102}

\begin{document}

\begin{tcolorbox}[colback=ucbblue!5!white, colframe=ucbblue, title=\textbf{Hoja de Registro y Comparación de Evidencia}]
    \centering \textbf{Trilogía de Validación: Cálculo Analítico $\rightarrow$ Simulación $\rightarrow$ Banco de Pruebas}
\end{tcolorbox}

\textbf{Equipo:} \hrulefill 

\vspace{0.2cm}
\noindent El Error Relativo debe calcularse respecto al valor teórico: 
\begin{equation*}
    \varepsilon_r = \left| \frac{\text{Medido} - \text{Teórico}}{\text{Teórico}} \right| \times 100\%
\end{equation*}

\renewcommand{\arraystretch}{1.8}
\noindent
\begin{tabularx}{\textwidth}{|>{\bfseries\RaggedRight}p{2.5cm}|>{\centering\arraybackslash}p{1.2cm}|>{\RaggedRight\arraybackslash}X|>{\RaggedRight\arraybackslash}X|>{\RaggedRight\arraybackslash}X|>{\RaggedRight\arraybackslash}X|>{\RaggedRight\arraybackslash}p{2.5cm}|}
\hline
\rowcolor{ucbblue!10}
\textbf{Magnitud} & \textbf{Var.} & \textbf{Teoría} & \textbf{SPICE} & \textbf{Medición} & \textbf{Error \%} & \textbf{Observación / Discrepancia} \\ \hline
Superposición Contribución 1 & $V_{a}^{(V1)}$ & $4.627\,\text{V}$ & & & & \\ \hline
Superposición Contribución 2 & $V_{a}^{(V2)}$ & $1.087\,\text{V}$ & & & & \\ \hline
Voltaje Total \newline (Nodal) & $V_{a}$ & $5.714\,\text{V}$ & & & & \\ \hline
Voltaje de Circ. Abierto & $V_{oc}$ & $5.714\,\text{V}$ & & & & \\ \hline
Resistencia de Thévenin & $R_{th}$ & $2.175\,\text{k}\Omega$ & & & & \\ \hline
Corriente de Norton & $I_{N}$ & $2.628\,\text{mA}$ & & \textit{(Indirecta)} & & \\ \hline
Carga 1 \newline ($R_L \approx 1\,\text{k}\Omega$) & $V_{L1}$ & $1.800\,\text{V}$ & & & & \\ \hline
Carga 2 \newline ($R_L \approx 2.2\,\text{k}\Omega$) & $V_{L2}$ & $2.874\,\text{V}$ & & & & \\ \hline
Carga 3 \newline ($R_L \approx 4.7\,\text{k}\Omega$) & $V_{L3}$ & $3.907\,\text{V}$ & & & & \\ \hline
Carga 4 \newline ($R_L \approx 10\,\text{k}\Omega$) & $V_{L4}$ & $4.694\,\text{V}$ & & & & \\ \hline
\end{tabularx}

\vspace{0.5cm}
\textbf{Conclusión de Discrepancias:} Identifique la fuente principal de error físico (no humano) observada en el banco. \\
\rule{\textwidth}{0.4pt} \\
\rule{\textwidth}{0.4pt}

\end{document}
